\chapter{扇区怎样组成 rootfs 的目录和文件}

ATA 只读写编号扇区，文件系统必须在扇区中约定超级块、inode、bitmap、目录项
和数据块的位置。文件变大以后，inode 还要通过多级索引寻找更多数据块。本章
从一个空镜像开始，创建目录和大文件，再重新挂载并运行检查器。

\chapterroute{分区、超级块和 inode $\rightarrow$ bitmap、目录项与缓存
$\rightarrow$ rootfs v2 盘面 $\rightarrow$ 多级块索引和稀疏文件
$\rightarrow$ 目录修改顺序 $\rightarrow$ 镜像生成、破坏和 fsck。}

\input{deepening/ch13-file-system}
\input{deepening/ch13-rootfs-v2}
